
\documentclass[UTF8,11pt,a4paper]{ctexart}
\usepackage[a4paper,left=23mm,right=23mm,top=24mm,bottom=24mm]{geometry}
\usepackage{fontspec}
\setmainfont{DejaVu Sans}
\setCJKmainfont{Noto Serif CJK SC}
\usepackage{amsmath,amssymb}
\usepackage{xcolor}
\usepackage[colorlinks=true,linkcolor=blue!45!black,unicode=true]{hyperref}
\usepackage{enumitem}
\usepackage{needspace}
\usepackage{fancyhdr}
\setlength{\parindent}{2em}
\setlength{\parskip}{0.25em}
\setlist[enumerate]{leftmargin=2.6em,itemsep=0.4em,topsep=0.5em}
\pagestyle{fancy}
\fancyhf{}
\fancyhead[L]{\small 三千弱水 · 大学化学高阶综合题库}
\fancyhead[R]{\small 待专家审校稿}
\fancyfoot[C]{\thepage}
\setlength{\headheight}{15pt}
\sloppy

\begin{document}
\begin{center}
{\LARGE\bfseries 大学化学高阶综合题 · 第16辑 (14题)}\\[0.8em]
{\large 题目、分步解析与数值核验}\\[0.5em]
{\small \textcolor{red!65!black}{审校稿：未经化学专家核准；禁止直接视为正式试卷或出版物}}
\end{center}
\vspace{1em}
\noindent 本册从仓库结构化 JSONL 自动生成。增补和修改应只编辑源题目数据，并重新执行统一编译。
\par 机械方程残差检查只能说明指定模型的内部数值一致，不保证模型现实可行性。第六辑另须完成跨辑模型去重审查。
\tableofcontents
\clearpage
\section*{第16辑 · 题目及解答}
\addcontentsline{toc}{section}{第16辑}
\subsection*{无机化学}
\addcontentsline{toc}{subsection}{第16辑 / 无机化学}
\Needspace{7\baselineskip}
\subsubsection*{UC-ADV16-I251 · Na₂O–SiO₂网络玻璃中非桥氧与Q³/Q⁴四面体}
\noindent\textbf{题目}\par
理想钠硅酸盐玻璃包含0.2000 mol Na₂O和0.8000 mol SiO₂，所有Si四配位，Na⁺仅用于电荷补偿，每个非桥氧携一个−1有效电荷。假设只有Q³和Q⁴四面体，没有Q²、三配位氧、铝杂质及其他电荷补偿机制。求总Si、Na、O摩尔数、非桥氧NBO摩尔数、每四面体非桥氧数、Q³与Q⁴摩尔数及分率、桥氧原子摩尔数、氧原子总量回代和每Si桥氧键平均数。计算桥氧时必须把两个四面体共用同一氧原子的事实纳入计数。\par
\medskip\noindent\textbf{分步解答}
\begin{enumerate}
\item 完整条件：理想钠硅酸盐玻璃包含0.2000 mol Na₂O和0.8000 mol SiO₂，所有Si四配位，Na⁺仅用于电荷补偿，每个非桥氧携一个−1有效电荷。假设只有Q³和Q⁴四面体，没有Q²、三配位氧、铝杂质及其他电荷补偿机制。求总Si、Na、O摩尔数、非桥氧NBO摩尔数、每四面体非桥氧数、Q³与Q⁴摩尔数及分率、桥氧原子摩尔数、氧原子总量回代和每Si桥氧键平均数。计算桥氧时必须把两个四面体共用同一氧原子的事实纳入计数。
\item 列出物料/电荷守恒与定量模型：nNBO=2nNa2O，nBO=(4nSi−nNBO)/2，nQ3=nNBO，nQ4=nSi−nQ3。
\item 明确测量单位、形式价态或标准态定义，并区分固定外库和封闭系统的计量。
\item 依据模型推导第一组变量以及在此范围采用的理想化假设。
\item 计算初级量：silicon\_mol = 0.80000000000000 mol；sodium\_mol = 0.40000000000000 mol；oxygen\_mol = 1.8000000000000 mol。
\item 由此得到中间结构与观测指标：NBO\_mol = 0.40000000000000 mol；NBO\_per\_tetra = 0.50000000000000 ratio；Q3\_mol = 0.40000000000000 mol。
\item 全部剩余数值和回代核对：Q4\_mol = 0.40000000000000 mol；Q3\_fraction = 0.50000000000000 fraction；Q4\_fraction = 0.50000000000000 fraction；BO\_mol = 1.4000000000000 mol；oxygen\_balance = 1.8000000000000 mol；bridging\_bonds\_per\_Si = 3.5000000000000 bonds/Si。
\item 把结果返回原始物料或校准方程，检验电荷、原子或浓度守恒，排除不可能的负浓度与超量回收。
\item 通过第二种代数表达式逐项复算数值向量，错误阈值1e−9相对精度。
\item 科学模型有效性、跨辑重复性及版权审核仍待专家审定，不进入可出版正式库。
\end{enumerate}
\noindent\textbf{答案摘要}\par
silicon\_mol = 0.8 mol；sodium\_mol = 0.4 mol；oxygen\_mol = 1.8 mol；NBO\_mol = 0.4 mol；NBO\_per\_tetra = 0.5 ratio；Q3\_mol = 0.4 mol；Q4\_mol = 0.4 mol；Q3\_fraction = 0.5 fraction；Q4\_fraction = 0.5 fraction；BO\_mol = 1.4 mol；oxygen\_balance = 1.8 mol；bridging\_bonds\_per\_Si = 3.5 bonds/Si\par
\noindent\textbf{模型前提：}
理想计量与仪器线性假设；第二路径数值复算不代表实验和原创性专家批准。\par
\noindent\textcolor{red!65!black}{[跨辑去重未完成 · 草稿]} \par
\medskip\hrule\medskip
\Needspace{7\baselineskip}
\subsubsection*{UC-ADV16-I252 · 重铬酸根与Fe²⁺酸性氧化的电子及原子守恒}
\noindent\textbf{题目}\par
在充分酸性的理想水溶液中定量完成Cr₂O₇²⁻+14H⁺+6Fe²⁺→2Cr³⁺+7H₂O+6Fe³⁺。初始Cr₂O₇²⁻=0.001000 mol、Fe²⁺=0.01000 mol，酸由外部池充足供给。没有其它氧化剂、还原剂和Cr配合反应，产物定量收集。求总电子转移mol及当量电量、Fe²⁺消耗和剩余、Fe³⁺与Cr³⁺生成量、酸消耗与水生成、Cr原子库存、Cr每个原子氧化态变化幅度及一摩尔重铬酸根的Fe²⁺需求数。说明虽然每个Cr(VI)接受3电子，一个重铬酸根含2个Cr共接受6电子。\par
\medskip\noindent\textbf{分步解答}
\begin{enumerate}
\item 完整条件：在充分酸性的理想水溶液中定量完成Cr₂O₇²⁻+14H⁺+6Fe²⁺→2Cr³⁺+7H₂O+6Fe³⁺。初始Cr₂O₇²⁻=0.001000 mol、Fe²⁺=0.01000 mol，酸由外部池充足供给。没有其它氧化剂、还原剂和Cr配合反应，产物定量收集。求总电子转移mol及当量电量、Fe²⁺消耗和剩余、Fe³⁺与Cr³⁺生成量、酸消耗与水生成、Cr原子库存、Cr每个原子氧化态变化幅度及一摩尔重铬酸根的Fe²⁺需求数。说明虽然每个Cr(VI)接受3电子，一个重铬酸根含2个Cr共接受6电子。
\item 列出物料/电荷守恒与定量模型：每mol Cr2O7²−接受6e−，FeII消耗6 mol、H+消耗14 mol、CrIII生成2 mol。
\item 明确测量单位、形式价态或标准态定义，并区分固定外库和封闭系统的计量。
\item 依据模型推导第一组变量以及在此范围采用的理想化假设。
\item 计算初级量：electron\_mol = 0.0060000000000000 mol；charge\_C = 578.91199272000 C；FeII\_consumed = 0.0060000000000000 mol。
\item 由此得到中间结构与观测指标：FeII\_left = 0.0040000000000000 mol；FeIII\_produced = 0.0060000000000000 mol；CrIII\_produced = 0.0020000000000000 mol。
\item 全部剩余数值和回代核对：acid\_consumed = 0.014000000000000 mol；water\_produced = 0.0070000000000000 mol；Cr\_atom\_mol = 0.0020000000000000 mol；oxidation\_drop\_per\_Cr = 3.0000000000000 oxidation；FeII\_per\_dichromate = 6.0000000000000 ratio。
\item 把结果返回原始物料或校准方程，检验电荷、原子或浓度守恒，排除不可能的负浓度与超量回收。
\item 通过第二种代数表达式逐项复算数值向量，错误阈值1e−9相对精度。
\item 科学模型有效性、跨辑重复性及版权审核仍待专家审定，不进入可出版正式库。
\end{enumerate}
\noindent\textbf{答案摘要}\par
electron\_mol = 0.006 mol；charge\_C = 578.911993 C；FeII\_consumed = 0.006 mol；FeII\_left = 0.004 mol；FeIII\_produced = 0.006 mol；CrIII\_produced = 0.002 mol；acid\_consumed = 0.014 mol；water\_produced = 0.007 mol；Cr\_atom\_mol = 0.002 mol；oxidation\_drop\_per\_Cr = 3 oxidation；FeII\_per\_dichromate = 6 ratio\par
\noindent\textbf{模型前提：}
理想计量与仪器线性假设；第二路径数值复算不代表实验和原创性专家批准。\par
\noindent\textcolor{red!65!black}{[跨辑去重未完成 · 草稿]} \par
\medskip\hrule\medskip
\Needspace{7\baselineskip}
\subsubsection*{UC-ADV16-I253 · Na₄Al₄Si₈O₂₄沸石交换Ca²⁺的理论容量}
\noindent\textbf{题目}\par
某假想无水沸石式量Na₄Al₄Si₈O₂₄中Na为+1、Al为+3、Si为+4、O为−2。每个AlO₄替换SiO₄所产生的一个负框架电荷由一个Na⁺补偿。Na、Al、Si、O的摩尔质量分别22.98977、26.98154、28.085、15.999 g/mol。假定所有Na交换位点可用，Ca²⁺离子交换时2个Ca²⁺恰取代4个Na⁺。求框架电荷、Na正电荷、盐式量质量、每g交换Na的mmol和meq、每式量可交换Ca的mol及每g交换Ca的mmol、Al/Si比、氧原子数、完全交换释放Na摩尔数和交换后电荷残差。不能把Ca²⁺的1 mol误当作1 meq。\par
\medskip\noindent\textbf{分步解答}
\begin{enumerate}
\item 完整条件：某假想无水沸石式量Na₄Al₄Si₈O₂₄中Na为+1、Al为+3、Si为+4、O为−2。每个AlO₄替换SiO₄所产生的一个负框架电荷由一个Na⁺补偿。Na、Al、Si、O的摩尔质量分别22.98977、26.98154、28.085、15.999 g/mol。假定所有Na交换位点可用，Ca²⁺离子交换时2个Ca²⁺恰取代4个Na⁺。求框架电荷、Na正电荷、盐式量质量、每g交换Na的mmol和meq、每式量可交换Ca的mol及每g交换Ca的mmol、Al/Si比、氧原子数、完全交换释放Na摩尔数和交换后电荷残差。不能把Ca²⁺的1 mol误当作1 meq。
\item 列出物料/电荷守恒与定量模型：摩尔质量M=4MNa+4MAl+8MSi+24MO；Na交换容量4×1000/M meq/g。
\item 明确测量单位、形式价态或标准态定义，并区分固定外库和封闭系统的计量。
\item 依据模型推导第一组变量以及在此范围采用的理想化假设。
\item 计算初级量：framework\_charge = -4.0000000000000 e；Na\_charge = 4.0000000000000 e；formula\_mass = 808.54124000000 g/mol。
\item 由此得到中间结构与观测指标：Na\_mmol\_g = 4.9471811728490 mmol/g；capacity\_meq\_g = 4.9471811728490 meq/g；Ca\_per\_formula = 2.0000000000000 mol。
\item 全部剩余数值和回代核对：Ca\_mmol\_g = 2.4735905864245 mmol/g；Al\_Si\_ratio = 0.50000000000000 ratio；oxygen\_count = 24.000000000000 atoms；Na\_released = 4.0000000000000 mol；charge\_balance = 0.0000000000000 e。
\item 把结果返回原始物料或校准方程，检验电荷、原子或浓度守恒，排除不可能的负浓度与超量回收。
\item 通过第二种代数表达式逐项复算数值向量，错误阈值1e−9相对精度。
\item 科学模型有效性、跨辑重复性及版权审核仍待专家审定，不进入可出版正式库。
\end{enumerate}
\noindent\textbf{答案摘要}\par
framework\_charge = -4 e；Na\_charge = 4 e；formula\_mass = 808.54124 g/mol；Na\_mmol\_g = 4.94718117 mmol/g；capacity\_meq\_g = 4.94718117 meq/g；Ca\_per\_formula = 2 mol；Ca\_mmol\_g = 2.47359059 mmol/g；Al\_Si\_ratio = 0.5 ratio；oxygen\_count = 24 atoms；Na\_released = 4 mol；charge\_balance = 0 e\par
\noindent\textbf{模型前提：}
理想计量与仪器线性假设；第二路径数值复算不代表实验和原创性专家批准。\par
\noindent\textcolor{red!65!black}{[跨辑去重未完成 · 草稿]} \par
\medskip\hrule\medskip
\Needspace{7\baselineskip}
\subsubsection*{UC-ADV16-I254 · 磷钨多金属氧酸盐Na₃[PW₁₂O₄₀]的电荷和质量分数}
\noindent\textbf{题目}\par
假设磷钨Keggin型杂多酸钠盐的完全脱水式量为Na₃[PW₁₂O₄₀]，所有Na按+1、P按+5、W按+6、O按−2形式价。摩尔质量分别为Na 22.98977、P 30.97376、W 183.84、O 15.999 g/mol。求复合阴离子电荷、12个W的形式总正电荷、40个氧总负电荷、盐式量阳离子加中心原子总正电荷与净电荷、W平均氧化态、复合阴离子和完整钠盐的摩尔质量、钨质量分数、氧质量分数、W/P原子数比和W(VI)形式d电子数。请强调形式氧化态是电荷清点，不等于真实分子轨道电荷分布。\par
\medskip\noindent\textbf{分步解答}
\begin{enumerate}
\item 完整条件：假设磷钨Keggin型杂多酸钠盐的完全脱水式量为Na₃[PW₁₂O₄₀]，所有Na按+1、P按+5、W按+6、O按−2形式价。摩尔质量分别为Na 22.98977、P 30.97376、W 183.84、O 15.999 g/mol。求复合阴离子电荷、12个W的形式总正电荷、40个氧总负电荷、盐式量阳离子加中心原子总正电荷与净电荷、W平均氧化态、复合阴离子和完整钠盐的摩尔质量、钨质量分数、氧质量分数、W/P原子数比和W(VI)形式d电子数。请强调形式氧化态是电荷清点，不等于真实分子轨道电荷分布。
\item 列出物料/电荷守恒与定量模型：Qanion=5+12×6−40×2=−3，M=3MNa+MP+12MW+40MO。
\item 明确测量单位、形式价态或标准态定义，并区分固定外库和封闭系统的计量。
\item 依据模型推导第一组变量以及在此范围采用的理想化假设。
\item 计算初级量：anion\_charge = -3.0000000000000 e；W\_charge\_total = 72.000000000000 e；oxygen\_charge\_total = -80.000000000000 e。
\item 由此得到中间结构与观测指标：salt\_positive\_charge = 80.000000000000 e；salt\_net\_charge = 0.0000000000000 e；W\_oxidation = 6.0000000000000 oxidation。
\item 全部剩余数值和回代核对：anion\_mass = 2877.0137600000 g/mol；salt\_mass = 2945.9830700000 g/mol；W\_fraction = 0.74884340730444 fraction；O\_fraction = 0.21723139094618 fraction；W\_P\_ratio = 12.000000000000 ratio；WVI\_d\_count = 0.0000000000000 electrons。
\item 把结果返回原始物料或校准方程，检验电荷、原子或浓度守恒，排除不可能的负浓度与超量回收。
\item 通过第二种代数表达式逐项复算数值向量，错误阈值1e−9相对精度。
\item 科学模型有效性、跨辑重复性及版权审核仍待专家审定，不进入可出版正式库。
\end{enumerate}
\noindent\textbf{答案摘要}\par
anion\_charge = -3 e；W\_charge\_total = 72 e；oxygen\_charge\_total = -80 e；salt\_positive\_charge = 80 e；salt\_net\_charge = 0 e；W\_oxidation = 6 oxidation；anion\_mass = 2877.01376 g/mol；salt\_mass = 2945.98307 g/mol；W\_fraction = 0.748843407 fraction；O\_fraction = 0.217231391 fraction；W\_P\_ratio = 12 ratio；WVI\_d\_count = 0 electrons\par
\noindent\textbf{模型前提：}
理想计量与仪器线性假设；第二路径数值复算不代表实验和原创性专家批准。\par
\noindent\textcolor{red!65!black}{[跨辑去重未完成 · 草稿]} \par
\medskip\hrule\medskip
\subsection*{分析化学}
\addcontentsline{toc}{subsection}{第16辑 / 分析化学}
\Needspace{7\baselineskip}
\subsubsection*{UC-ADV16-A251 · 酶促UV反应的吸光度时间斜率与摩尔初速}
\noindent\textbf{题目}\par
某酶分析比色皿光程l=1.000 cm，产物P在340 nm处的摩尔吸光系数ε=6220 L/(mol cm)，初始30.00 s内实测扣除空白吸光度从0.1000线性增加到0.2866，分析混合体积V=3.000 mL。假设仅P在此波长吸收，产物按1:1来自底物且在30秒内初速保持恒定。求吸光度增量、每秒吸光度斜率、产物浓度增加、浓度变化速率mol/(L s)、体积内产物生成mol/s和30s生成总量、产物浓度在30s时值、初始时产物浓度、相对初始吸光增加倍率及酶活性单位U（定义1U=1 μmol/min）。不要把mL体积直接当作L。\par
\medskip\noindent\textbf{分步解答}
\begin{enumerate}
\item 完整条件：某酶分析比色皿光程l=1.000 cm，产物P在340 nm处的摩尔吸光系数ε=6220 L/(mol cm)，初始30.00 s内实测扣除空白吸光度从0.1000线性增加到0.2866，分析混合体积V=3.000 mL。假设仅P在此波长吸收，产物按1:1来自底物且在30秒内初速保持恒定。求吸光度增量、每秒吸光度斜率、产物浓度增加、浓度变化速率mol/(L s)、体积内产物生成mol/s和30s生成总量、产物浓度在30s时值、初始时产物浓度、相对初始吸光增加倍率及酶活性单位U（定义1U=1 μmol/min）。不要把mL体积直接当作L。
\item 列出物料/电荷守恒与定量模型：Δc=ΔA/(εl)，dc/dt=ΔA/(εl Δt)，enzyme\_U=(dc/dt)×V(L)×60×1e6。
\item 明确测量单位、形式价态或标准态定义，并区分固定外库和封闭系统的计量。
\item 依据模型推导第一组变量以及在此范围采用的理想化假设。
\item 计算初级量：delta\_absorbance = 0.18660000000000 AU；absorbance\_slope = 0.0062200000000000 AU/s；delta\_product\_M = 0.000030000000000000 mol/L。
\item 由此得到中间结构与观测指标：rate\_Ms = 0.0000010000000000000 mol/(Ls)；rate\_mol\_s = 3.0000000000000e-9 mol/s；moles\_product\_30s = 9.0000000000000e-8 mol。
\item 全部剩余数值和回代核对：product\_conc\_at30\_M = 0.000046077170418006 mol/L；product\_conc\_at0\_M = 0.000016077170418006 mol/L；absorbance\_ratio\_final\_initial = 2.8660000000000 ratio；enzyme\_activity\_U = 0.18000000000000 umol/min。
\item 把结果返回原始物料或校准方程，检验电荷、原子或浓度守恒，排除不可能的负浓度与超量回收。
\item 通过第二种代数表达式逐项复算数值向量，错误阈值1e−9相对精度。
\item 科学模型有效性、跨辑重复性及版权审核仍待专家审定，不进入可出版正式库。
\end{enumerate}
\noindent\textbf{答案摘要}\par
delta\_absorbance = 0.1866 AU；absorbance\_slope = 0.00622 AU/s；delta\_product\_M = 3e-05 mol/L；rate\_Ms = 1e-06 mol/(Ls)；rate\_mol\_s = 3e-09 mol/s；moles\_product\_30s = 9e-08 mol；product\_conc\_at30\_M = 4.60771704e-05 mol/L；product\_conc\_at0\_M = 1.60771704e-05 mol/L；absorbance\_ratio\_final\_initial = 2.866 ratio；enzyme\_activity\_U = 0.18 umol/min\par
\noindent\textbf{模型前提：}
理想计量与仪器线性假设；第二路径数值复算不代表实验和原创性专家批准。\par
\noindent\textcolor{red!65!black}{[跨辑去重未完成 · 草稿]} \par
\medskip\hrule\medskip
\Needspace{7\baselineskip}
\subsubsection*{UC-ADV16-A252 · 色谱内标法双峰响应因子与稀释回算}
\noindent\textbf{题目}\par
某GC分析组分X时采用同样量纲的内标Y。标准液中CX,std=10.00 mg/L、CY,std=5.000 mg/L，响应峰面积分别AX,std=8000、AY,std=4000 counts；样品被5.000倍定容稀释，并加入同浓度Y内标，测得AX,unk=4800、AY,unk=3000 counts。假设响应比AX/AY=k×CX/CY，标准液与稀释样的内标Y浓度相同，仪器漂移对两峰同比例。求标准响应比、标准浓度比、相对响应因子k、未知样峰比、未知稀释后X浓度、原样X浓度、未知样Y信号漂移分数、校正后未知X峰面积、样品相对标准的X浓度比及如果错误只用X绝对峰面积得到的原样浓度。须区分校准响应因子与样品稀释倍数。\par
\medskip\noindent\textbf{分步解答}
\begin{enumerate}
\item 完整条件：某GC分析组分X时采用同样量纲的内标Y。标准液中CX,std=10.00 mg/L、CY,std=5.000 mg/L，响应峰面积分别AX,std=8000、AY,std=4000 counts；样品被5.000倍定容稀释，并加入同浓度Y内标，测得AX,unk=4800、AY,unk=3000 counts。假设响应比AX/AY=k×CX/CY，标准液与稀释样的内标Y浓度相同，仪器漂移对两峰同比例。求标准响应比、标准浓度比、相对响应因子k、未知样峰比、未知稀释后X浓度、原样X浓度、未知样Y信号漂移分数、校正后未知X峰面积、样品相对标准的X浓度比及如果错误只用X绝对峰面积得到的原样浓度。须区分校准响应因子与样品稀释倍数。
\item 列出物料/电荷守恒与定量模型：k=(AX/AY)std/(CX/CY)std，Cunk=(AX/AY)unk/k×CY，Coriginal=5Cunk。
\item 明确测量单位、形式价态或标准态定义，并区分固定外库和封闭系统的计量。
\item 依据模型推导第一组变量以及在此范围采用的理想化假设。
\item 计算初级量：std\_signal\_ratio = 2.0000000000000 ratio；std\_concentration\_ratio = 2.0000000000000 ratio；relative\_response\_factor = 1.0000000000000 ratio。
\item 由此得到中间结构与观测指标：unknown\_signal\_ratio = 1.6000000000000 ratio；diluted\_X\_mgL = 8.0000000000000 mg/L；original\_X\_mgL = 40.000000000000 mg/L。
\item 全部剩余数值和回代核对：internal\_signal\_drift = 0.75000000000000 fraction；corrected\_unknown\_X\_area = 6400.0000000000 counts；diluted\_to\_std\_X\_ratio = 0.80000000000000 ratio；uncorrected\_original\_X\_mgL = 30.000000000000 mg/L。
\item 把结果返回原始物料或校准方程，检验电荷、原子或浓度守恒，排除不可能的负浓度与超量回收。
\item 通过第二种代数表达式逐项复算数值向量，错误阈值1e−9相对精度。
\item 科学模型有效性、跨辑重复性及版权审核仍待专家审定，不进入可出版正式库。
\end{enumerate}
\noindent\textbf{答案摘要}\par
std\_signal\_ratio = 2 ratio；std\_concentration\_ratio = 2 ratio；relative\_response\_factor = 1 ratio；unknown\_signal\_ratio = 1.6 ratio；diluted\_X\_mgL = 8 mg/L；original\_X\_mgL = 40 mg/L；internal\_signal\_drift = 0.75 fraction；corrected\_unknown\_X\_area = 6400 counts；diluted\_to\_std\_X\_ratio = 0.8 ratio；uncorrected\_original\_X\_mgL = 30 mg/L\par
\noindent\textbf{模型前提：}
理想计量与仪器线性假设；第二路径数值复算不代表实验和原创性专家批准。\par
\noindent\textcolor{red!65!black}{[跨辑去重未完成 · 草稿]} \par
\medskip\hrule\medskip
\Needspace{7\baselineskip}
\subsubsection*{UC-ADV16-A253 · ICP电沉积二价金属的库仑电量与回收质量}
\noindent\textbf{题目}\par
某电化学痕量分析对Zn²⁺进行恒电流沉积Zn²⁺+2e⁻→Zn(s)，电流I=0.01000 A、电解时间t=600.0 s，沉积法拉第效率η=0.9500，F=96485.33212 C/mol、MZn=65.38 g/mol。样品总体积100.0 mL且Zn²⁺初始浓度0.001000 mol/L，保证沉积量不超过库存。求总电量、有效电量、有效电子mol、沉积Zn的mol与质量mg、对应样品浓度降低量(mol/L)、沉积回收占初始Zn摩尔数比例、非金属沉积消耗电量、若误以单电子沉积计算的Zn质量mg及由此得到的质量高估倍数。题设忽略锌表面再溶解、氧析出以外其他副反应。\par
\medskip\noindent\textbf{分步解答}
\begin{enumerate}
\item 完整条件：某电化学痕量分析对Zn²⁺进行恒电流沉积Zn²⁺+2e⁻→Zn(s)，电流I=0.01000 A、电解时间t=600.0 s，沉积法拉第效率η=0.9500，F=96485.33212 C/mol、MZn=65.38 g/mol。样品总体积100.0 mL且Zn²⁺初始浓度0.001000 mol/L，保证沉积量不超过库存。求总电量、有效电量、有效电子mol、沉积Zn的mol与质量mg、对应样品浓度降低量(mol/L)、沉积回收占初始Zn摩尔数比例、非金属沉积消耗电量、若误以单电子沉积计算的Zn质量mg及由此得到的质量高估倍数。题设忽略锌表面再溶解、氧析出以外其他副反应。
\item 列出物料/电荷守恒与定量模型：Q=It，Qeff=ηQ，nZn=Qeff/(2F)，Δc=nZn/V。
\item 明确测量单位、形式价态或标准态定义，并区分固定外库和封闭系统的计量。
\item 依据模型推导第一组变量以及在此范围采用的理想化假设。
\item 计算初级量：total\_charge\_C = 6.0000000000000 C；effective\_charge\_C = 5.7000000000000 C；electron\_mol = 0.000059076337042721 mol。
\item 由此得到中间结构与观测指标：deposited\_Zn\_mol = 0.000029538168521361 mol；deposited\_Zn\_mg = 1.9312054579266 mg；concentration\_drop\_M = 0.00029538168521361 mol/L。
\item 全部剩余数值和回代核对：fraction\_sample\_Zn\_recovered = 0.29538168521361 fraction；parasitic\_charge\_C = 0.30000000000000 C；wrong\_single\_electron\_mass\_mg = 3.8624109158531 mg；overestimate\_factor = 2.0000000000000 ratio。
\item 把结果返回原始物料或校准方程，检验电荷、原子或浓度守恒，排除不可能的负浓度与超量回收。
\item 通过第二种代数表达式逐项复算数值向量，错误阈值1e−9相对精度。
\item 科学模型有效性、跨辑重复性及版权审核仍待专家审定，不进入可出版正式库。
\end{enumerate}
\noindent\textbf{答案摘要}\par
total\_charge\_C = 6 C；effective\_charge\_C = 5.7 C；electron\_mol = 5.9076337e-05 mol；deposited\_Zn\_mol = 2.95381685e-05 mol；deposited\_Zn\_mg = 1.93120546 mg；concentration\_drop\_M = 0.000295381685 mol/L；fraction\_sample\_Zn\_recovered = 0.295381685 fraction；parasitic\_charge\_C = 0.3 C；wrong\_single\_electron\_mass\_mg = 3.86241092 mg；overestimate\_factor = 2 ratio\par
\noindent\textbf{模型前提：}
理想计量与仪器线性假设；第二路径数值复算不代表实验和原创性专家批准。\par
\noindent\textcolor{red!65!black}{[跨辑去重未完成 · 草稿]} \par
\medskip\hrule\medskip
\Needspace{7\baselineskip}
\subsubsection*{UC-ADV16-A254 · 离子色谱加标回收率及原样稀释因子的计量校正}
\noindent\textbf{题目}\par
某离子色谱方法测定样品阴离子X，先将原样稀释10.00倍，未加标的稀释样测得仪器表观浓度C0=1.800 mg/L。向同体积的另一份稀释样加入最终浓度贡献2.000 mg/L的X标准溶液，且加标定容保持两份终体积相等，测得Csp=3.700 mg/L。假设仪器响应对内源X与标准X具有相同的恒定回收系数r，且无额外背景，因此r=(Csp−C0)/2.000。求绝对加标增量、回收率及百分数、经回收校正的稀释样真实X浓度、原样真实浓度、若忽略回收所得原样浓度、回收校正浓度的相对增幅、标准加标实测增量相对理论比例、校正后的加标样真实总浓度及经校正计算的原样每1 L X质量mg。不能将加标后未经回收校正的表观值直接当原样真值。\par
\medskip\noindent\textbf{分步解答}
\begin{enumerate}
\item 完整条件：某离子色谱方法测定样品阴离子X，先将原样稀释10.00倍，未加标的稀释样测得仪器表观浓度C0=1.800 mg/L。向同体积的另一份稀释样加入最终浓度贡献2.000 mg/L的X标准溶液，且加标定容保持两份终体积相等，测得Csp=3.700 mg/L。假设仪器响应对内源X与标准X具有相同的恒定回收系数r，且无额外背景，因此r=(Csp−C0)/2.000。求绝对加标增量、回收率及百分数、经回收校正的稀释样真实X浓度、原样真实浓度、若忽略回收所得原样浓度、回收校正浓度的相对增幅、标准加标实测增量相对理论比例、校正后的加标样真实总浓度及经校正计算的原样每1 L X质量mg。不能将加标后未经回收校正的表观值直接当原样真值。
\item 列出物料/电荷守恒与定量模型：r=(Csp−C0)/Cadded；Creal,dil=C0/r，Creal,orig=10Creal,dil。
\item 明确测量单位、形式价态或标准态定义，并区分固定外库和封闭系统的计量。
\item 依据模型推导第一组变量以及在此范围采用的理想化假设。
\item 计算初级量：measured\_spike\_increment\_mgL = 1.9000000000000 mg/L；recovery\_fraction = 0.95000000000000 fraction；recovery\_percent = 95.000000000000 percent。
\item 由此得到中间结构与观测指标：corrected\_diluted\_conc = 1.8947368421053 mg/L；corrected\_original\_conc = 18.947368421053 mg/L；uncorrected\_original\_conc = 18.000000000000 mg/L。
\item 全部剩余数值和回代核对：relative\_correction\_increase = 0.052631578947368 fraction；measured\_theoretical\_increment\_ratio = 0.95000000000000 ratio；corrected\_spiked\_total\_conc = 3.8947368421053 mg/L；original\_mg\_per\_L = 18.947368421053 mg。
\item 把结果返回原始物料或校准方程，检验电荷、原子或浓度守恒，排除不可能的负浓度与超量回收。
\item 通过第二种代数表达式逐项复算数值向量，错误阈值1e−9相对精度。
\item 科学模型有效性、跨辑重复性及版权审核仍待专家审定，不进入可出版正式库。
\end{enumerate}
\noindent\textbf{答案摘要}\par
measured\_spike\_increment\_mgL = 1.9 mg/L；recovery\_fraction = 0.95 fraction；recovery\_percent = 95 percent；corrected\_diluted\_conc = 1.89473684 mg/L；corrected\_original\_conc = 18.9473684 mg/L；uncorrected\_original\_conc = 18 mg/L；relative\_correction\_increase = 0.0526315789 fraction；measured\_theoretical\_increment\_ratio = 0.95 ratio；corrected\_spiked\_total\_conc = 3.89473684 mg/L；original\_mg\_per\_L = 18.9473684 mg\par
\noindent\textbf{模型前提：}
理想计量与仪器线性假设；第二路径数值复算不代表实验和原创性专家批准。\par
\noindent\textcolor{red!65!black}{[跨辑去重未完成 · 草稿]} \par
\medskip\hrule\medskip
\subsection*{有机化学}
\addcontentsline{toc}{subsection}{第16辑 / 有机化学}
\Needspace{7\baselineskip}
\subsubsection*{UC-ADV16-O248 · 芳香亲电取代邻位/对位统计简并度和产物分流}
\noindent\textbf{题目}\par
单取代苯A在理想硝化条件下仅发生单次邻位或对位取代，单个邻位碳上的拟一级微观速度常数ko=0.02000 s⁻¹，单个对位碳kp=0.03000 s⁻¹。由于有两个对称邻位和一个对位，邻位总通道速率2ko，对位通道kp，硝化剂浓度由外池保持不变。初始A=0.01000 mol/L，反应时间30.00 s；没有间位、二次硝化和其它副反应。求两通道及总拟一级速度、邻/对产物分率、30秒残余A和两种产物浓度、产物浓度比、总转化率、半衰期及单个邻位相对于对位的微观速度比。解释宏观区域选择性必须计算位点简并度。\par
\medskip\noindent\textbf{分步解答}
\begin{enumerate}
\item 题目及全部条件：单取代苯A在理想硝化条件下仅发生单次邻位或对位取代，单个邻位碳上的拟一级微观速度常数ko=0.02000 s⁻¹，单个对位碳kp=0.03000 s⁻¹。由于有两个对称邻位和一个对位，邻位总通道速率2ko，对位通道kp，硝化剂浓度由外池保持不变。初始A=0.01000 mol/L，反应时间30.00 s；没有间位、二次硝化和其它副反应。求两通道及总拟一级速度、邻/对产物分率、30秒残余A和两种产物浓度、产物浓度比、总转化率、半衰期及单个邻位相对于对位的微观速度比。解释宏观区域选择性必须计算位点简并度。
\item 建立机理或分配方程：kobs=2ko+kp，f\_o=2ko/kobs，Po=A0(1−e\textasciicircum{}(−kobst))f\_o。
\item 明确化学骨架和官能团守恒关系，区分并行通道、独立反应与物理萃取。
\item 求解概率与反应计量前，将体积和速率常数换为统一单位。
\item 前四项：ortho\_rate = 0.04000000000000 s\textasciicircum{}-1；para\_rate = 0.03000000000000 s\textasciicircum{}-1；total\_rate = 0.07000000000000 s\textasciicircum{}-1；ortho\_fraction = 0.5714285714286 fraction。
\item 中间结果：para\_fraction = 0.4285714285714 fraction；remaining\_A = 0.001224564282530 mol/L；ortho\_product = 0.005014534695697 mol/L；para\_product = 0.003760901021773 mol/L。
\item 余下数值答案：ortho\_para\_ratio = 1.333333333333 ratio；conversion = 0.8775435717470 fraction；half\_life = 9.902102579428 s；single\_site\_rate\_ratio = 0.6666666666667 ratio。
\item 检查物质的量总库存、通道分率及凝胶分支参数的逻辑限值。
\item 由独立代数表达式重算各数值，采用不高于1e−9相对误差标准。
\item 本题尚待专家科学和跨辑原创性评审，不属于已批准正式出版题。
\end{enumerate}
\noindent\textbf{答案摘要}\par
ortho\_rate = 0.04 s\textasciicircum{}-1；para\_rate = 0.03 s\textasciicircum{}-1；total\_rate = 0.07 s\textasciicircum{}-1；ortho\_fraction = 0.571428571 fraction；para\_fraction = 0.428571429 fraction；remaining\_A = 0.00122456428 mol/L；ortho\_product = 0.0050145347 mol/L；para\_product = 0.00376090102 mol/L；ortho\_para\_ratio = 1.33333333 ratio；conversion = 0.877543572 fraction；half\_life = 9.90210258 s；single\_site\_rate\_ratio = 0.666666667 ratio\par
\noindent\textbf{模型前提：}
全部模型均为题设理想情况；专家学术和原创性审核待完成。\par
\noindent\textcolor{red!65!black}{[跨辑去重未完成 · 草稿]} \par
\medskip\hrule\medskip
\Needspace{7\baselineskip}
\subsubsection*{UC-ADV16-O249 · A₃/B₂三官能–二官能步增长聚合的凝胶临界转化率}
\noindent\textbf{题目}\par
某理想随机步增长聚合中，A₃分子每个带3个A型官能团，B₂分子每个带2个B型官能团，只允许A–B反应，且没有环化。起始加入A₃ 2.000 mol与B₂ 3.000 mol，双方可反应官能团均为6.000 mol，各自转化率为相同p。采用Flory–Stockmayer分支参数α=(3−1)(2−1)p²，凝胶判据α≥1。比较p=0.6000和0.8000，求临界p、两状态α、p=0.8时已形成的A–B键mol与A/B各自剩余的官能团mol、两个状态是否超过临界的标志、p=0.8相对临界的转化率差、初始总官能团数、剩余总官能团数和初态A/B官能团数量比。不能凭凝胶判据直接推算胶团质量分数。\par
\medskip\noindent\textbf{分步解答}
\begin{enumerate}
\item 题目及全部条件：某理想随机步增长聚合中，A₃分子每个带3个A型官能团，B₂分子每个带2个B型官能团，只允许A–B反应，且没有环化。起始加入A₃ 2.000 mol与B₂ 3.000 mol，双方可反应官能团均为6.000 mol，各自转化率为相同p。采用Flory–Stockmayer分支参数α=(3−1)(2−1)p²，凝胶判据α≥1。比较p=0.6000和0.8000，求临界p、两状态α、p=0.8时已形成的A–B键mol与A/B各自剩余的官能团mol、两个状态是否超过临界的标志、p=0.8相对临界的转化率差、初始总官能团数、剩余总官能团数和初态A/B官能团数量比。不能凭凝胶判据直接推算胶团质量分数。
\item 建立机理或分配方程：pc=1/√2；α=2p²；生成AB键量=6p，剩余A和B各6(1−p)。
\item 明确化学骨架和官能团守恒关系，区分并行通道、独立反应与物理萃取。
\item 求解概率与反应计量前，将体积和速率常数换为统一单位。
\item 前四项：gel\_critical\_p = 0.7071067811865 fraction；branching\_p0p6 = 0.7200000000000 ratio；branching\_p0p8 = 1.280000000000 ratio；formed\_bonds\_mol = 4.800000000000 mol。
\item 中间结果：A\_groups\_left\_mol = 1.200000000000 mol；B\_groups\_left\_mol = 1.200000000000 mol；gel\_flag\_p0p6 = 0.000000000000 boolean；gel\_flag\_p0p8 = 1.000000000000 boolean。
\item 余下数值答案：p\_above\_threshold = 0.09289321881345 fraction；initial\_total\_groups\_mol = 12.00000000000 mol；remaining\_total\_groups\_mol = 2.400000000000 mol；initial\_A\_B\_group\_ratio = 1.000000000000 ratio。
\item 检查物质的量总库存、通道分率及凝胶分支参数的逻辑限值。
\item 由独立代数表达式重算各数值，采用不高于1e−9相对误差标准。
\item 本题尚待专家科学和跨辑原创性评审，不属于已批准正式出版题。
\end{enumerate}
\noindent\textbf{答案摘要}\par
gel\_critical\_p = 0.707106781 fraction；branching\_p0p6 = 0.72 ratio；branching\_p0p8 = 1.28 ratio；formed\_bonds\_mol = 4.8 mol；A\_groups\_left\_mol = 1.2 mol；B\_groups\_left\_mol = 1.2 mol；gel\_flag\_p0p6 = 0 boolean；gel\_flag\_p0p8 = 1 boolean；p\_above\_threshold = 0.0928932188 fraction；initial\_total\_groups\_mol = 12 mol；remaining\_total\_groups\_mol = 2.4 mol；initial\_A\_B\_group\_ratio = 1 ratio\par
\noindent\textbf{模型前提：}
全部模型均为题设理想情况；专家学术和原创性审核待完成。\par
\noindent\textcolor{red!65!black}{[跨辑去重未完成 · 草稿]} \par
\medskip\hrule\medskip
\Needspace{7\baselineskip}
\subsubsection*{UC-ADV16-O250 · 一次与三次分液萃取的有机物分配回收比较}
\noindent\textbf{题目}\par
中性有机化合物X在水相与有机溶剂之间的平衡分配系数D=[X]org/[X]aq=5.000，水相体积100.0 mL，初始X物质的量0.01000 mol。假定两相不互溶、每次接触都达到平衡且不电离或分解。比较用50.00 mL有机溶剂一次萃取，以及将同样50.00 mL溶剂等分为三份、每次用新鲜溶剂连续萃取，萃取后只保留水层进入下次。求单次的水相残留分率、回收分率及水相剩余mol，三次方案每步残留因子、最终水相分率与剩余mol、总回收分率、相对单次回收率提高幅度、两方案剩余摩尔数比、单次方案水相最终浓度。须说明“多次小体积”的回收优势来自每次更换新鲜有机相。\par
\medskip\noindent\textbf{分步解答}
\begin{enumerate}
\item 题目及全部条件：中性有机化合物X在水相与有机溶剂之间的平衡分配系数D=[X]org/[X]aq=5.000，水相体积100.0 mL，初始X物质的量0.01000 mol。假定两相不互溶、每次接触都达到平衡且不电离或分解。比较用50.00 mL有机溶剂一次萃取，以及将同样50.00 mL溶剂等分为三份、每次用新鲜溶剂连续萃取，萃取后只保留水层进入下次。求单次的水相残留分率、回收分率及水相剩余mol，三次方案每步残留因子、最终水相分率与剩余mol、总回收分率、相对单次回收率提高幅度、两方案剩余摩尔数比、单次方案水相最终浓度。须说明“多次小体积”的回收优势来自每次更换新鲜有机相。
\item 建立机理或分配方程：f=Vaq/(Vaq+D Vorg)，单次f1=100/(100+5×50)，三次f3=[100/(100+5×50/3)]³。
\item 明确化学骨架和官能团守恒关系，区分并行通道、独立反应与物理萃取。
\item 求解概率与反应计量前，将体积和速率常数换为统一单位。
\item 前四项：single\_remain\_fraction = 0.2857142857143 fraction；single\_recovered\_fraction = 0.7142857142857 fraction；single\_remain\_mol = 0.002857142857143 mol；per\_step\_remain\_fraction = 0.5454545454545 fraction。
\item 中间结果：triple\_remain\_fraction = 0.1622839969947 fraction；triple\_remain\_mol = 0.001622839969947 mol；triple\_recovered\_fraction = 0.8377160030053 fraction；improvement\_fraction = 0.1234302887195 fraction。
\item 余下数值答案：remaining\_ratio\_triple\_single = 0.5679939894816 ratio；single\_aqueous\_concentration = 0.02857142857143 mol/L。
\item 检查物质的量总库存、通道分率及凝胶分支参数的逻辑限值。
\item 由独立代数表达式重算各数值，采用不高于1e−9相对误差标准。
\item 本题尚待专家科学和跨辑原创性评审，不属于已批准正式出版题。
\end{enumerate}
\noindent\textbf{答案摘要}\par
single\_remain\_fraction = 0.285714286 fraction；single\_recovered\_fraction = 0.714285714 fraction；single\_remain\_mol = 0.00285714286 mol；per\_step\_remain\_fraction = 0.545454545 fraction；triple\_remain\_fraction = 0.162283997 fraction；triple\_remain\_mol = 0.00162283997 mol；triple\_recovered\_fraction = 0.837716003 fraction；improvement\_fraction = 0.123430289 fraction；remaining\_ratio\_triple\_single = 0.567993989 ratio；single\_aqueous\_concentration = 0.0285714286 mol/L\par
\noindent\textbf{模型前提：}
全部模型均为题设理想情况；专家学术和原创性审核待完成。\par
\noindent\textcolor{red!65!black}{[跨辑去重未完成 · 草稿]} \par
\medskip\hrule\medskip
\subsection*{物理化学}
\addcontentsline{toc}{subsection}{第16辑 / 物理化学}
\Needspace{7\baselineskip}
\subsubsection*{UC-ADV16-P248 · Morse非谐振子的谱项、跃迁波数与解离能}
\noindent\textbf{题目}\par
一种假想双原子分子的振动项采用Morse模型G(v)=ωe(v+1/2)−ωexe(v+1/2)²，单位cm⁻¹，给定谐振常数ωe=2000.0 cm⁻¹，非谐参数xe=0.01000，因此ωexe=20.00 cm⁻¹。忽略转动、电子激发态和同位素效应。求v=0、1、2的振动谱项、0→1及1→2跃迁波数、两跃迁间隔差、0→1光谱波长nm、势阱底至解离极限的Morse理论解离能De/hc=ωe²/(4ωexe)、基态零点能及其相对于De的比例、0→2跃迁波数。以λ(nm)=10⁷/ν(cm⁻¹)换算；零点能不是零，解离能De也不能混作从v0开始的D0。\par
\medskip\noindent\textbf{分步解答}
\begin{enumerate}
\item 理想物理化学实验模型：一种假想双原子分子的振动项采用Morse模型G(v)=ωe(v+1/2)−ωexe(v+1/2)²，单位cm⁻¹，给定谐振常数ωe=2000.0 cm⁻¹，非谐参数xe=0.01000，因此ωexe=20.00 cm⁻¹。忽略转动、电子激发态和同位素效应。求v=0、1、2的振动谱项、0→1及1→2跃迁波数、两跃迁间隔差、0→1光谱波长nm、势阱底至解离极限的Morse理论解离能De/hc=ωe²/(4ωexe)、基态零点能及其相对于De的比例、0→2跃迁波数。以λ(nm)=10⁷/ν(cm⁻¹)换算；零点能不是零，解离能De也不能混作从v0开始的D0。
\item 推导热力学或统计方程：G(v)=2000(v+½)−20(v+½)²，De/hc=2000²/(4×20)。
\item 将频率、nm与m以及mol/m³的单位换算明确列入公式，避免无量纲参数与原始实测量混淆。
\item 由独立基础量建立最小完整方程组，并检查模型适用的压力、温度及几何边界。
\item 主要数值：G0 = 995.00000000000 cm\textasciicircum{}-1；G1 = 2955.0000000000 cm\textasciicircum{}-1；G2 = 4875.0000000000 cm\textasciicircum{}-1；transition01 = 1960.0000000000 cm\textasciicircum{}-1。
\item 后续可观测指标：transition12 = 1920.0000000000 cm\textasciicircum{}-1；transition\_difference = 40.000000000000 cm\textasciicircum{}-1；wavelength01\_nm = 5102.0408163265 nm；well\_depth\_cmInv = 50000.000000000 cm\textasciicircum{}-1。
\item 其余输出及回代：zero\_point\_cmInv = 995.00000000000 cm\textasciicircum{}-1；zero\_point\_fraction = 0.019900000000000 fraction；transition02 = 3880.0000000000 cm\textasciicircum{}-1。
\item 逐项核查谱线能差、物质的量通量或统计分率归一，检验单位量纲和物理范围。
\item 以另一种代数表达式独立计算并对照全部数值向量，允许不超过1e−9相对偏差。
\item 模型结果尚未经专家科学性及跨辑原创性审核，维持待审稿不可正式发布状态。
\end{enumerate}
\noindent\textbf{答案摘要}\par
G0 = 995 cm\textasciicircum{}-1；G1 = 2955 cm\textasciicircum{}-1；G2 = 4875 cm\textasciicircum{}-1；transition01 = 1960 cm\textasciicircum{}-1；transition12 = 1920 cm\textasciicircum{}-1；transition\_difference = 40 cm\textasciicircum{}-1；wavelength01\_nm = 5102.04082 nm；well\_depth\_cmInv = 50000 cm\textasciicircum{}-1；zero\_point\_cmInv = 995 cm\textasciicircum{}-1；zero\_point\_fraction = 0.0199 fraction；transition02 = 3880 cm\textasciicircum{}-1\par
\noindent\textbf{模型前提：}
经典统计与输运近似仅限于题设条件；专家审核未完成。\par
\noindent\textcolor{red!65!black}{[跨辑去重未完成 · 草稿]} \par
\medskip\hrule\medskip
\Needspace{7\baselineskip}
\subsubsection*{UC-ADV16-P249 · 纳米孔内N₂ Knudsen扩散速率与孔径的线性关系}
\noindent\textbf{题目}\par
298.15 K，某直通圆柱孔内N₂处于分子–孔壁碰撞控制的Knudsen极限，摩尔质量M=0.0280134 kg/mol，R=8.314462618。孔半径r1=10.00 nm、r2=20.00 nm，孔长L=0.001000 m，两端浓度差Δc=1.000 mol/m³。Knudsen扩散系数DK=(2/3)r√[8RT/(πM)]，稳态扩散通量J=DK Δc/L。求两个孔的DK与J、两扩散系数比及通量比、平均热速率尺度、如果有效通道总横截面积为0.01000 m²时r1对应的总摩尔流率与分子数流率、两个孔的一维特征扩散时间τ=L²/(2DK)。忽略孔道曲折、孔隙率校正和连续介质的气体分子间碰撞。\par
\medskip\noindent\textbf{分步解答}
\begin{enumerate}
\item 理想物理化学实验模型：298.15 K，某直通圆柱孔内N₂处于分子–孔壁碰撞控制的Knudsen极限，摩尔质量M=0.0280134 kg/mol，R=8.314462618。孔半径r1=10.00 nm、r2=20.00 nm，孔长L=0.001000 m，两端浓度差Δc=1.000 mol/m³。Knudsen扩散系数DK=(2/3)r√[8RT/(πM)]，稳态扩散通量J=DK Δc/L。求两个孔的DK与J、两扩散系数比及通量比、平均热速率尺度、如果有效通道总横截面积为0.01000 m²时r1对应的总摩尔流率与分子数流率、两个孔的一维特征扩散时间τ=L²/(2DK)。忽略孔道曲折、孔隙率校正和连续介质的气体分子间碰撞。
\item 推导热力学或统计方程：D\_K=(2r/3)√(8RT/πM)，J=DK/L，τ=L²/(2DK)。
\item 将频率、nm与m以及mol/m³的单位换算明确列入公式，避免无量纲参数与原始实测量混淆。
\item 由独立基础量建立最小完整方程组，并检查模型适用的压力、温度及几何边界。
\item 主要数值：D1 = 0.0000031646842538393 m2/s；D2 = 0.0000063293685076786 m2/s；flux1 = 0.0031646842538393 mol/(m2s)；flux2 = 0.0063293685076786 mol/(m2s)。
\item 后续可观测指标：diffusion\_ratio = 2.0000000000000 ratio；flux\_ratio = 2.0000000000000 ratio；thermal\_speed = 474.70263807590 m/s；molar\_flow1 = 0.000031646842538393 mol/s。
\item 其余输出及回代：molecules\_per\_second1 = 1.9058174037576e+19 molecules/s；transit\_time1 = 0.15799364482994 s；transit\_time2 = 0.078996822414971 s。
\item 逐项核查谱线能差、物质的量通量或统计分率归一，检验单位量纲和物理范围。
\item 以另一种代数表达式独立计算并对照全部数值向量，允许不超过1e−9相对偏差。
\item 模型结果尚未经专家科学性及跨辑原创性审核，维持待审稿不可正式发布状态。
\end{enumerate}
\noindent\textbf{答案摘要}\par
D1 = 3.16468425e-06 m2/s；D2 = 6.32936851e-06 m2/s；flux1 = 0.00316468425 mol/(m2s)；flux2 = 0.00632936851 mol/(m2s)；diffusion\_ratio = 2 ratio；flux\_ratio = 2 ratio；thermal\_speed = 474.702638 m/s；molar\_flow1 = 3.16468425e-05 mol/s；molecules\_per\_second1 = 1.9058174e+19 molecules/s；transit\_time1 = 0.157993645 s；transit\_time2 = 0.0789968224 s\par
\noindent\textbf{模型前提：}
经典统计与输运近似仅限于题设条件；专家审核未完成。\par
\noindent\textcolor{red!65!black}{[跨辑去重未完成 · 草稿]} \par
\medskip\hrule\medskip
\Needspace{7\baselineskip}
\subsubsection*{UC-ADV16-P250 · 经验Hill结合曲线的配体半响应与对数斜率}
\noindent\textbf{题目}\par
某理想蛋白质配体结合信号在本题工作区间由经验Hill公式θ=L\textasciicircum{}n/(Kd\textasciicircum{}n+L\textasciicircum{}n)刻画，Hill系数n=2.000、半饱和自由配体浓度Kd=0.01000 mol/L。外部配体储液池维持L=0.005000、0.01000、0.02000 mol/L三个状态，忽略配体耗竭和真实位点异质性。求三个状态结合分率与未结合分率、三个结合/未结合的几率比、在L=Kd时的对数灵敏度dθ/dlnL=nθ(1−θ)及对L的导数、由Kd/2升到2Kd的覆盖分率增量、两极端状态结合几率比的比值以及Hill对数几率曲线的斜率。Hill系数是宏观协同性，不直接证明有两个独立物理结合位点。\par
\medskip\noindent\textbf{分步解答}
\begin{enumerate}
\item 理想物理化学实验模型：某理想蛋白质配体结合信号在本题工作区间由经验Hill公式θ=L\textasciicircum{}n/(Kd\textasciicircum{}n+L\textasciicircum{}n)刻画，Hill系数n=2.000、半饱和自由配体浓度Kd=0.01000 mol/L。外部配体储液池维持L=0.005000、0.01000、0.02000 mol/L三个状态，忽略配体耗竭和真实位点异质性。求三个状态结合分率与未结合分率、三个结合/未结合的几率比、在L=Kd时的对数灵敏度dθ/dlnL=nθ(1−θ)及对L的导数、由Kd/2升到2Kd的覆盖分率增量、两极端状态结合几率比的比值以及Hill对数几率曲线的斜率。Hill系数是宏观协同性，不直接证明有两个独立物理结合位点。
\item 推导热力学或统计方程：θ=L²/(Kd²+L²)，odds=θ/(1−θ)，dθ/dlnL=2θ(1−θ)。
\item 将频率、nm与m以及mol/m³的单位换算明确列入公式，避免无量纲参数与原始实测量混淆。
\item 由独立基础量建立最小完整方程组，并检查模型适用的压力、温度及几何边界。
\item 主要数值：theta\_half = 0.20000000000000 fraction；theta\_equal = 0.50000000000000 fraction；theta\_double = 0.80000000000000 fraction；unbound\_half = 0.80000000000000 fraction。
\item 后续可观测指标：unbound\_equal = 0.50000000000000 fraction；unbound\_double = 0.20000000000000 fraction；odds\_half = 0.25000000000000 ratio；odds\_equal = 1.0000000000000 ratio。
\item 其余输出及回代：odds\_double = 4.0000000000000 ratio；log\_sensitivity\_at\_equal = 0.50000000000000 fraction；derivative\_dtheta\_dL\_equal = 50.000000000000 M\textasciicircum{}-1；coverage\_increment = 0.60000000000000 fraction；extreme\_odds\_ratio = 16.000000000000 ratio；Hill\_slope = 2.0000000000000 ratio。
\item 逐项核查谱线能差、物质的量通量或统计分率归一，检验单位量纲和物理范围。
\item 以另一种代数表达式独立计算并对照全部数值向量，允许不超过1e−9相对偏差。
\item 模型结果尚未经专家科学性及跨辑原创性审核，维持待审稿不可正式发布状态。
\end{enumerate}
\noindent\textbf{答案摘要}\par
theta\_half = 0.2 fraction；theta\_equal = 0.5 fraction；theta\_double = 0.8 fraction；unbound\_half = 0.8 fraction；unbound\_equal = 0.5 fraction；unbound\_double = 0.2 fraction；odds\_half = 0.25 ratio；odds\_equal = 1 ratio；odds\_double = 4 ratio；log\_sensitivity\_at\_equal = 0.5 fraction；derivative\_dtheta\_dL\_equal = 50 M\textasciicircum{}-1；coverage\_increment = 0.6 fraction；extreme\_odds\_ratio = 16 ratio；Hill\_slope = 2 ratio\par
\noindent\textbf{模型前提：}
经典统计与输运近似仅限于题设条件；专家审核未完成。\par
\noindent\textcolor{red!65!black}{[跨辑去重未完成 · 草稿]} \par
\medskip\hrule\medskip
\end{document}
